% Faithful English typesetting of the frozen 7 October 2026 release basis.
% Input hashes and editorial changes are recorded in TYPESETTING_LEDGER.md.
\documentclass[11pt,letterpaper]{article}
\usepackage[T1]{fontenc}
\usepackage[utf8]{inputenc}
\usepackage{lmodern}
\usepackage[margin=1in]{geometry}
\usepackage{amsmath,amssymb,amsthm,mathtools}
\usepackage{microtype}
\usepackage{xurl}
\usepackage[hidelinks]{hyperref}
\hypersetup{pdftitle={A polynomial-dimensional refinement of sharp simplex stability},
 pdfauthor={AI-assisted research supplement to entry005},
 pdfsubject={Faithful typeset derivative of the frozen 7 October 2026 mathematical release basis}}
\pdfinfoomitdate=1
\pdftrailerid{}
\newtheorem*{maintheorem}{Theorem}
\newcommand{\EE}{\mathbb E}
\newcommand{\PP}{\mathbb P}
\newcommand{\RR}{\mathbb R}
\newcommand{\BB}{B_2^d}
\newcommand{\conv}{\operatorname{conv}}
\newcommand{\supp}{\operatorname{supp}}
\newcommand{\diam}{\operatorname{diam}}
\newcommand{\per}{\operatorname{per}}
\newcommand{\Rproj}{R_{\mathrm{proj}}}
\newcommand{\pos}[1]{(#1)_+}
\setlength{\emergencystretch}{2em}
\allowdisplaybreaks[1]
\title{A polynomial-dimensional refinement\\of sharp simplex stability}
\author{AI-assisted research supplement to entry005}
\date{7 October 2026}
\begin{document}
\maketitle
\vspace{-1.5em}
\begin{abstract}
For the projection-cone invariant of the pinned source manuscripts, the
sharp exponent $1/(d-1)$ in the centroid containment estimate for every
maximum-volume inscribed simplex admits an explicit coefficient bounded by
$2^{20}d^6$. The proof uses first absolute determinant moments, a weighted
selection of support anchors, the actual convex-body cone law, projection
caps, and retention of the prescribed maximum simplex. An explicit
square-pyramid construction gives the necessary coefficient bound
$(d+1)[d(d+1)]^{1/(d-1)}$. This is a typeset derivative of the frozen
written refinement and its companion lower-bound proof; it makes no
first-discovery, best-known-result, human peer-review, or proof-assistant
formalization claim.
\end{abstract}

\section*{Result and conventions}
A convex body is a compact convex set with nonempty interior. We use exactly
the invariant and excess of the pinned manuscripts. The projection body
$\Pi K$ is defined by
\[
 h_{\Pi K}(u)=|K\mid u^\perp|_{d-1}\qquad (\|u\|=1),
\]
where the right-hand side is the $(d-1)$-dimensional volume of the orthogonal
projection. The operator $\mathcal P$ denotes a pyramid over $K$, with the
convention in entry005 v3, Section~1 \cite{entry}. In particular,
$\mathcal P K$ has dimension $d+1$; the projection ratio is evaluated in
the dimension of its argument. These are the source definitions, not a
different projection-body invariant. Put
\begin{align*}
 \Rproj(K)&=\frac{|\Pi K|}{|K|^{d-1}},\\
 a(K)&=\left(\frac d{d+1}\right)^d
       \frac{\Rproj(\mathcal P K)}{\Rproj(K)}-1,\\
 e(K)&=a(K)-\frac1{d+1},\\
 E(K,S)&=\inf\{t\ge0:K\subset z_S+(1+t)(S-z_S)\},
\end{align*}
where $z_S$ is the centroid of $S$.

\begin{maintheorem}
For every $d\ge3$, every full-dimensional convex body $K\subset\RR^d$,
and every maximum-volume inscribed simplex $S$,
\[
 E(K,S)\le G_d e(K)^{1/(d-1)}
          \le 2^{20}d^6 e(K)^{1/(d-1)}.
\]
The exponent is unchanged and optimal, by the actual simplex-truncation
family in the pinned manuscript \cite{sharp,truncation}.
\end{maintheorem}

The exact constants are
\begin{align*}
 m&=d-1, & R&=d\sqrt{d+2}, & M&=4dR,\\
 C&=(d+1)^2(d+2), & L&=2\sqrt m\,(R+1),\\
 J&=\frac d2(2R)^d\bigl[1+M+d2^{d-1}M\bigr],\\
 G_d&=8(R+1)dL(JC)^{1/m},
 & e_0&=\frac1{JC(8dL)^m}.
\end{align*}
For $0\le e\le e_0$, the stronger local estimate with coefficient $G_d/2$
holds. Moreover,
\[
 G_d e_0^{1/m}=R+1,
\]
which makes the global branch transparent.

The main structural strengthening is valid for probability laws without
any convex-body realization assumption: a centered full-dimensional
probability law of bounded support admits an assignment to $d+1$ affinely
independent support points with mean distance at most
\[
 \diam(\supp\nu)\,\frac{(d+1)(d+2)}2\,\frac{B-A}{B}.
\]
Distances and diameter may be taken in any fixed norm. This assertion
requires no dispersion assumption, inverse determinant, covariance
estimate, second determinant moment, or small-probability conditioning loss.
The refinement is relative to the supplied argument; no first-discovery
or best-known-result claim is made.

\section{The stronger integrated assignment theorem}
Let $\nu$ be centered with bounded support spanning $\RR^d$. Let
$X_1,X_2,\ldots$ be independent with law $\nu$, and define the
\emph{first absolute determinant moments}
\begin{align*}
 A&=\EE\,|\det(X_1,\ldots,X_d)|,\\
 B&=\EE\,|\det((X_0,1),\ldots,(X_d,1))|, & D&=B-A.
\end{align*}
Full-dimensionality and centering imply $A>0$ and $B>0$. Set $n=d+1$ and
$N=n(n+1)/2$. For an ordered $d$-tuple $\mathbf x$, define
\begin{align*}
 F_{\mathbf x}(y)&=\det((x_1,1),\ldots,(x_d,1),(y,1)),\\
 P(\mathbf x)&=\EE\pos{F_{\mathbf x}(X)}, &
 N_-(\mathbf x)&=\EE\pos{-F_{\mathbf x}(X)}.
\end{align*}
Centering gives $D=2\EE\min(P,N_-)\ge0$. For independent test points $Y,Z$
define the nonnegative witness
\begin{align*}
 \psi_{\mathbf x}(Y,Z)
 &=\min\{\pos{F_{\mathbf x}(Y)},\pos{-F_{\mathbf x}(Z)}\}\\
 &\quad+\min\{\pos{-F_{\mathbf x}(Y)},\pos{F_{\mathbf x}(Z)}\}.
\end{align*}
Its expectation over all $d+2$ independent samples is at most $D$.

For anchors $W=(w_0,\ldots,w_d)$, let $V(W)$ be the absolute lifted
determinant, \emph{not} the simplex volume $V(W)/d!$. Define $H(W)$ by
integrating in $x$ the sum of these $N$ witnesses:
\begin{enumerate}
\item For every $i$, the base is the $d$ other anchors and the tests are
      $w_i,x$.
\item For every $i<j$, the base is $x$ and the $d-1$ anchors other than
      $i,j$, and the tests are $w_i,w_j$.
\end{enumerate}
Order bases by a fixed convention; $\psi$ is unaffected by changing the
overall determinant sign. Each summand uses $d+2$ independent samples
before conditioning, so $\EE H\le ND$.

\subsection*{Selection without an expectation-of-ratio error}
On anchor tuples define the probability measure
\[
 d\PP_V(W)=\frac{V(W)}B\,d\nu^{\otimes(d+1)}(W).
\]
This is a probability because $\EE V=B$. It gives full mass to support
tuples with $V>0$. Define $H/V$ only there. Then
\[
 \EE_{\PP_V}(H/V)=\frac{\EE[H\mathbf1_{\{V>0\}}]}B\le\frac{ND}B.
\]
Consequently some support tuple with $V>0$ satisfies $H/V\le ND/B$.
Indeed, if the integrable ratio were strictly larger everywhere on a
full-measure set, its expectation would be strictly larger too. No compact
minimum on the open set $\{V>0\}$ is assumed. When $D=0$, the ratio is zero
$\PP_V$-almost surely.

\subsection*{Pointwise clipped-witness bound}
Fix such a tuple and let $\alpha_i(x)$ be its affine barycentric coordinates.
Let $r(x)$ be the first index attaining the largest positive coefficient;
some positive coefficient exists because $\sum_i\alpha_i=1$. This is Borel
measurable. Put
\[
 \varphi(x)=\sum_i\min\{1,\pos{-\alpha_i(x)}\}
       +\sum_{i<j}\min\{\pos{\alpha_i(x)},\pos{\alpha_j(x)}\}.
\]
The first witnesses equal $V\min\{1,\pos{-\alpha_i}\}$. When
$\alpha_i,\alpha_j>0$, the pair witness is $V\min(\alpha_i,\alpha_j)$,
because the two tested determinants have opposite signs. When either is
nonpositive, the displayed lower bound is zero. Thus $V\EE\varphi\le H$.

Let $\Delta=\diam(\supp\nu)$ in this subsection. If some $\alpha_i\le-1$,
then $\varphi\ge1$ and
$\|x-w_r\|\le\Delta\le\Delta\varphi$. Otherwise every negative magnitude
is below $1$, so clipping is inactive. Let
\[
 N_x=\sum_i\pos{-\alpha_i},\qquad U_x=\sum_{i\ne r}\pos{\alpha_i}.
\]
The pairs containing $r$ give
$U_x\le\sum_{i<j}\min\{\pos{\alpha_i},\pos{\alpha_j}\}$. Also
\[
 x-w_r=\sum_{i\ne r}\alpha_i(x)(w_i-w_r),
\]
so $\|x-w_r\|\le\Delta(N_x+U_x)\le\Delta\varphi$. Therefore, in all cases,
\[
 h:=\EE\|X-w_{r(X)}\|
 \le\Delta\EE\varphi\le\Delta H/V\le\Delta ND/B.
\]
This is the entire stronger assignment theorem. In the Euclidean unit
ball, $\Delta\le2$, giving $h\le(d+1)(d+2)D/B$.

\subsection*{Quantitative origin inclusion when dispersion is supplied}
If $\EE\pos{u\cdot X}\ge\eta$ for every unit $u$ and $h\le\eta-b$ with
$b>0$, then the assigned law has $\EE\pos{u\cdot w_r}\ge b$. At least one
anchor has $u\cdot w_i\ge b$ in every direction. Thus
$b\BB\subset T=\conv(w_i)$. This implication is separate from the
assignment theorem and is exactly where cone-body geometry will enter.

\section[Prescribed-simplex normalization and two elementary improvements]{Prescribed-simplex normalization and\\two elementary improvements}
Fix the originally prescribed maximum simplex $S$; do not select a
replacement maximum. Apply an invertible affine map sending $S$ to the
regular simplex $\Delta$ centered at zero with inradius $1$. Its vertices
$v_i$ have $\|v_i\|=d$ and $v_i\cdot v_j=-d$ for $i\ne j$.

For every $x\in K$, replacing vertex $i$ by $x$ changes the simplex volume
by the factor $|\alpha_i(x)|$. Maximality gives $|\alpha_i(x)|\le1$. Since
$\sum_i\alpha_i=1$,
\begin{align*}
 \|x\|^2&=d\left[(d+1)\sum_i\alpha_i(x)^2-1\right]\\
 &\le d[(d+1)^2-1]=d^2(d+2).
\end{align*}
Therefore
\[
 \BB\subset\Delta\subset K\subset R\BB,\qquad R=d\sqrt{d+2}.
\]
Also, dilation of $\Delta$ by $1+t$ about its centroid is equivalent to
$\alpha_i\ge-t/(d+1)$. The same $|\alpha_i|\le1$ immediately yields the
global bound
\[
 E(K,\Delta)\le d+1.
\]
This is stronger than using $K\subset R\Delta$ for the global branch.

\section{Actual cone-law dispersion without arbitrary-law substitution}
For this normalized convex body, define $\nu$ as the pushforward of
$h_K(u)dS_K(u)/(d|K|)$ under $u\mapsto u/h_K(u)$. The classical identities
$\int h_K\,dS_K=d|K|$ and $\int u\,dS_K=0$ make it a centered probability
law, supported on $\partial K^\circ\subset\BB$. Its support spans $\RR^d$
because every orthogonal projection of $K$ has positive $(d-1)$-volume.
The source cone representation \cite[Proposition~4.1]{entry} gives
\[
 a(K)=\frac{B}{(d+1)A},\qquad
 \frac DB=\frac{(d+1)e}{1+(d+1)e}.
\]
Thus Section~1 gives anchors and a measurable assignment satisfying
\[
 h\le\frac{Ce}{1+(d+1)e}\le Ce,\qquad C=(d+1)^2(d+2).
\]
Cauchy's projection formula, including its factor $1/2$, gives
\[
 \EE|u\cdot X|=\frac{2\pi_K(u)}{d|K|},\qquad
 \EE\pos{u\cdot X}=\frac{\pi_K(u)}{d|K|}
\]
for unit $u$, where $\pi_K(u)=|K\mid u^\perp|_{d-1}$ and centering gives
the second equality.

The geometric lower bound is particularly simple. Slice $K$ by lines
parallel to $u$ over $K\mid u^\perp$. Every nonempty chord has length at most
the directional width $w_K(u)=h_K(u)+h_K(-u)$. Fubini therefore gives
\[
 |K|\le w_K(u)\pi_K(u)\le2R\pi_K(u).
\]
It follows that
\[
 \EE\pos{u\cdot X}\ge\frac1{2dR}=\frac2M,\qquad M=4dR.
\]
This is valid for actual convex-body cone laws. It is not asserted for
arbitrary centered laws merely because their support lies in a ball.

If $h\le1/M$, Section~1 implies $M^{-1}\BB\subset T\subset\BB$. Hence
$P=T^\circ$ is a genuine enclosing simplex with
\[
 \BB\subset K\subset P\subset M\BB.
\]
Each selected anchor lies on $\partial K^\circ$, so
$h_K(w_i)=h_P(w_i)=1$.

\section{Complete assignment-to-projection conversion}
Let $q_i$ be the vertex of $P$ opposite the facet with polar normal $w_i$,
and let $\lambda_i>0$ be the barycentric coordinates of zero in $T$. The
cone law of $P$ is $\sum_i\lambda_i\delta_{w_i}$: it is supported on exactly
these facet polar images, is centered, and has unique barycentric weights.

Let $p_i=\PP(r(X)=i)$ and $c=\sum_i p_iw_i$. Centering gives $\|c\|\le h$.
The barycentric functions of $T$ are
\[
 \alpha_i(x)=\lambda_i(1-q_i\cdot x).
\]
Thus $p_i-\lambda_i=-\lambda_i q_i\cdot c$ and
$\sum_i|p_i-\lambda_i|\le Mh$. Comparing the assigned law and $\nu_P$,
for every unit $u$,
\[
 \left|\EE_\nu|u\cdot X|-\EE_{\nu_P}|u\cdot X|\right|\le(1+M)h,
\]
and consequently
\[
 \left|\frac{\pi_K(u)}{|K|}-\frac{\pi_P(u)}{|P|}\right|
 \le\frac d2(1+M)h.
\]
Since $h_P$ is $M$-Lipschitz and $h_P(w_i)=1$,
\[
 0\le z:=\int(h_P(x)-1)\,d\nu(x)\le Mh\le1.
\]
By homogeneity in $x=u/h_K(u)$, the integral equals
$V(K[d-1],P)/|K|-1$. Minkowski's first inequality, equivalently the first
variation of Brunn--Minkowski, yields
\[
 \frac{|P|}{|K|}\le(1+z)^d,\qquad
 |P|-|K|\le d2^{d-1}Mh\,|K|.
\]
Because $\supp\nu_P\subset\BB$, we have $\pi_P(u)/|P|\le d/2$. Combining
these facts and $|K|\le(2R)^d$ proves
\[
 0\le\pi_P(u)-\pi_K(u)\le Jh
\]
with $J$ as stated. Nonnegativity uses the genuine inclusion $K\subset P$.

\section{An improved projection cap using the radius of K}
Suppose $\BB\subset K\subset P$, $K\subset R\BB$, and all projection
deficits are at most $\varepsilon$. Let $s=d_H(K,P)$. The conclusion is
\[
 s\le2\sqrt{d-1}(R+1)\varepsilon^{1/(d-1)}
\]
whenever $\sqrt{d-1}\,\varepsilon^{1/(d-1)}\le1/2$. No bound on the outer
radius of $P$ is needed for this statement, apart from compactness.

For $s>0$, choose a farthest $q\in P$, its nearest $k\in K$, and
$n=(q-k)/s$. The projection optimality condition gives $n\cdot k=h_K(n)$.
Crucially,
\[
 \|q\|\le R+s.
\]
Choose a unit $u\perp n$ and project onto $u^\perp$, of dimension $m=d-1$.
Let $\tau=s/[2(R+s+1)]<1$. The projected $P$ contains the ball
$(1-\tau)\overline q+\tau B_2^m$. Its lowest $n$-coordinate is at least
\[
 n\cdot q-\tau(\|q\|+1)\ge h_K(n)+s/2,
\]
so the ball is disjoint from projected $K$. It contains a cube of side
$2\tau/\sqrt m$, hence
\[
 \varepsilon\ge\left[\frac{s}{\sqrt m\,(R+s+1)}\right]^m.
\]
Putting $t=\sqrt m\,\varepsilon^{1/m}$, we have $s\le t(R+s+1)$.
If $t\le1/2$, this rearranges to $s\le2t(R+1)$. The case $\varepsilon=0$
follows from the same cap inequality. This proves the claim.

This step avoids transporting the large polar radius $M$ into every
geometric factor. Hereafter $L=2\sqrt m\,(R+1)$.

\section{Local gate and retention of the original maximum simplex}
Assume $0\le e\le e_0$. Since $J\ge M$ and $8dL\ge1$,
\[
 h\le Ce\le\frac1{J(8dL)^m}\le\frac1M.
\]
Section~4 applies with $\varepsilon=Jh$. Write $y=(JCe)^{1/m}$. Then
\[
 y\le\frac1{8dL},\qquad
 \sqrt m\,y\le\frac1{16d(R+1)}\le\frac12.
\]
The improved cap gives
\[
 s\le L(Jh)^{1/m}\le Ly\le\frac1{8d}.
\]
In particular $P\subset(R+s)\BB\subset(R+1)\BB$. Also
$h_K\ge h_P-s\ge(1-s)h_P$, because $h_P\ge1$, so $(1-s)P\subset K$.

The fixed original $\Delta$ is maximum. Therefore
\[
 \frac{|\Delta|}{|P|}\ge(1-s)^d\ge1-ds.
\]
Put $\delta=ds\le1/8$. Express the vertices of $\Delta$ in barycentric
coordinates in $P$, producing a nonnegative column-stochastic
$(d+1)\times(d+1)$ matrix $W$ with
$|\det W|=|\Delta|/|P|\ge1-\delta$. Hadamard gives each column norm at
least $1-\delta$, and its largest coordinate is at least the squared norm,
hence at least $1-2\delta$.

These dominant coordinates occupy distinct rows: if two occupied the same
row, the independent-row-selection interpretation of the permanent would
give
\[
 |\det W|\le\per(W)\le1-(1-2\delta)^2\le4\delta<1-\delta,
\]
a contradiction. Match vertices using these distinct rows. Each vertex
of $\Delta$ differs from its matched vertex of $P$ by at most
\[
 2\delta\diam(P)\le4(R+1)ds.
\]
Every vertex of $P$ is matched. Thus
\[
 K\subset P\subset\Delta+4(R+1)ds\,\BB
       \subset[1+4(R+1)ds]\Delta,
\]
using $\BB\subset\Delta$. This proves
\[
 E(K,\Delta)\le4(R+1)dL(JC)^{1/m}e^{1/m}
                =\frac{G_d}{2}e^{1/m}.
\]
All inequalities include $e=0$: then $h=0$, projection deficits vanish,
$s=0$, and the same matching identifies the simplices.

\section{Global branch and original centroid}
By the exact definitions,
\[
 G_d e_0^{1/m}=R+1\ge d+1.
\]
For $e>e_0$, the universal maximality bound $E(K,\Delta)\le d+1$ therefore
gives $E(K,\Delta)\le G_d e^{1/m}$. For $e\le e_0$, Section~6 applies.
This closes the global theorem without any enormous minimum threshold
entering the final constant.

The affine normalization used the prescribed $S$ throughout. An invertible
affine map sends its centroid to the centroid of $\Delta$ and intertwines
homotheties about these centroids. Pulling back the final containment gives
exactly the stated dilation about $z_S$. The source invariant $a$, hence
$e$, is affine invariant \cite{entry,released}. No assertion that an
original Euclidean unit ball is preserved under affine pullback is needed.

\section{Explicit polynomial bound for the coefficient}
For $d\ge3$,
\begin{gather*}
 C\le3d^3,\qquad M\ge1,\qquad
 1+M+d2^{d-1}M\le d2^dM,\\
 J\le d^2M(4R)^d,\qquad R+1\le2R,\qquad L\le4\sqrt d\,R.
\end{gather*}
Consequently
\begin{align*}
 G_d&\le256d^{3/2}R^3[12d^5MR]^{1/(d-1)}\\
    &=256d^{3/2}R^3[48d^6R^2]^{1/(d-1)}.
\end{align*}
Since $R^2=d^2(d+2)\le(5/3)d^3$,
\[
 G_d\le256(5/3)^{3/2}d^6[80d^9]^{1/(d-1)}.
\]
For integer $d\ge3$, $d\le3^{(d-1)/2}$: this follows by induction, since
$(d+1)/d\le4/3<\sqrt3$. Therefore
\[
 [80d^9]^{1/(d-1)}\le\sqrt{80}\,3^{9/2}<1280.
\]
Also $(5/3)^{3/2}<9/4$. Hence
\[
 G_d<576\cdot1280\,d^6=737280d^6<2^{20}d^6.
\]
This polynomial is deliberately nonoptimized. The exact displayed $G_d$
is substantially smaller in many dimensions.

\section{Sharpness and remaining scope}
The exponent cannot improve: the pinned equal-intercept simplex
truncations have
\[
 E=(d+1)t,\qquad
 e\sim\frac{d(d-1)}{(d+1)^2}t^{d-1}
 \quad(t\downarrow0)
\]
for the maximum simplex used in the obstruction \cite{sharp,truncation}.
The direct actual-polytope construction proved in Appendix~\ref{app:square}
gives the stronger global coefficient obstruction
\[
 G_d\ge(d+1)[d(d+1)]^{1/(d-1)}.
\]
Its order is $d+1+2\log d+O((\log d)^2/d)$, leaving a gap between linear
and $d^6$ growth. It rules out a dimension-free coefficient, while the
theorem above rules out the necessity of superpolynomial dimension
dependence for this invariant and every-maximum-simplex metric.

The abstract assignment theorem controls mean transportation distance,
not support Hausdorff distance. Rare atoms can remain far away while
carrying negligible mass. The cone identities, centering, enclosing
inclusion, first variation, and projection caps are essential to the
geometric conversion. The companion \nolinkurl{weighted_anchors.md} in
the frozen release basis gives the detailed arbitrary-law theorem and
exact stress tests.

The frozen refinement records an independent analytic model review. Such
review is not human peer review or Lean formalization; no priority claim
or result about a different projection-body invariant is made. The
cone-law identity and affine invariance are imported from the pinned
source; the classical Cauchy, surface-area, first-variation, and
Brunn--Minkowski identities are explicitly identified inputs.

The release basis intentionally omits separate anisotropic and
simultaneous-truncation research. No theorem from that material is needed
here or endorsed by this release basis.

A focused next target is $G_d=O(d)$, which remains open here. The exact
$G_d$ above is asymptotic to $64d^6$. The remaining loss lies in Euclidean
cap and retention geometry. A relative cap estimate intrinsic to the
selected polar simplex could be investigated using the general-norm
assignment theorem; such an improvement is not proved here.

\section{Sources and literature positioning}
The exact source bytes and public commit URLs are recorded in the
accompanying source-pin file. The original sharp manuscript is
\cite{sharp}; entry005 v3 \cite{entry} and the released explicit modulus
\cite{released} supply the invariant's cone representation, its affine
invariance, and the older stability argument. This argument changes anchor
selection and dimension estimates, not the definition of the invariant.

The frozen refinement records a primary-literature inspection on
7 October 2026 with the following scope:
\begin{itemize}
\item B\"or\"oczky's \emph{The stability of the Rogers--Shephard inequality
and of some related inequalities} \cite{boroczky} has a Theorem~3 concerning
$|K|^{d-1}|\Pi^*K|$, the polar projection-body volume, and a Banach--Mazur
distance with a selectable simplex. Its stated upper modulus is
$d^{88d}\varepsilon^{1/d}$. This is a different deficit from $e(K)$; no
comparison identifying them has been proved here. Therefore the present
refinement is not asserted to improve that theorem.
\item B\"or\"oczky and Henk's \emph{Cone-volume measure and stability}
\cite{bh} concerns cone-volume concentration and the $U$-functional,
rather than the first absolute determinant difference used here. It does
not by itself give a realization theorem for arbitrary centered laws on
prescribed polar support points.
\item B\"or\"oczky, Fodor and Hug's \emph{Strengthened inequalities for the
mean width and the $\ell$-norm} \cite{bfh} studies different extremal
functionals and John/L\"owner positioning. It provides context, not an
imported estimate for $e(K)$.
\end{itemize}
The targeted primary-source search recorded in the frozen refinement did
not establish a prior theorem identical to the weighted first-moment
assignment theorem or the polynomial constant above. This is not an
exhaustive literature search and does not justify a novelty claim.

\appendix
\section{A stronger realized lower bound for the endpoint coefficient}
\label{app:square}
For every $d\ge3$, a universal coefficient in the every-maximum-simplex
estimate must satisfy
\[
 G_d\ge(d+1)[d(d+1)]^{1/(d-1)}.
\]
This is realized by the explicit full-dimensional polytope and maximum
simplex below. The frozen companion proof records an independent analytic
audit of the construction, maximality, barycentric excess, invariant
calculation, and direct facet certificate. No arbitrary law is treated as
a convex-body realization. The generic facet formula is from the pinned
entry005 v3, Section~1 \cite{entry}: if $u_i$ is facet area times outward
unit normal and $b_i$ is facet area times support number, then
\[
 H=\sum_{|I|=d}|\det(u_i:i\in I)|,\qquad
 L=\sum_{|I|=d+1}|\det((u_i,b_i):i\in I)|,
\]
\[
 a(K)=\frac{L}{d|K|H}.
\]
Here $H$ and $L$ are facet-minor sums local to this appendix, not the
assignment witness $H$ or cap constant $L$ of the main proof. The optional
pyramid calculation provides a second route; the direct certificate avoids
needing that recursion as an extra input.

\subsection{An elementary global bound on the geometric excess}
If $S=\conv(v_0,\ldots,v_d)$ is maximum-volume in $K$, its barycentric
coordinates satisfy
\[
 |\alpha_i(x)|\le1\qquad(x\in K).
\]
Indeed, replacing $v_i$ by $x$ multiplies volume by $|\alpha_i(x)|$.
Centroid dilation has the exact characterization
\[
 x\in z_S+(1+t)(S-z_S)
 \quad\Longleftrightarrow\quad
 \alpha_i(x)\ge-\frac{t}{d+1}\quad\text{for every }i.
\]
Consequently
\begin{equation}\label{eq:appendix-global}
 E(K,S)\le d+1.
\end{equation}
This applies to every maximum simplex without a normalization or a
small-deficit hypothesis. It is much sharper than the coarse outer-radius
bound used as the large-deficit branch in the supplied proof.

\subsection{The square-pyramid construction}
Take
\[
 C_d=\conv(0,e_1,e_2,e_1+e_2,e_3,\ldots,e_d),\qquad
 S_d=\conv(0,e_1,\ldots,e_d).
\]
This is the standard unit square followed by $d-2$ pyramids. The square
has $a=1/2$, and each pyramid adds one to $1/a$, so
\begin{equation}\label{eq:appendix-invariant}
 a(C_d)=\frac1d,\qquad e(C_d)=\frac1{d(d+1)}.
\end{equation}

To check maximality directly, a full-dimensional vertex simplex must
contain every $e_3,\ldots,e_d$, since no other vertex has a nonzero
corresponding coordinate. It must also contain exactly three square
vertices. Every such triangle has area $1/2$, so every full-dimensional
vertex simplex has volume $1/d!$. The absolute determinant is separately
convex in its vertices, hence no simplex with arbitrary vertices in $C_d$
has larger volume. In particular, $S_d$ is maximum.

The point $e_1+e_2$ has barycentric coordinates
\[
 (-1,1,1,0,\ldots,0)
\]
in $S_d$. Hence $E(C_d,S_d)\ge d+1$, with equality by
\eqref{eq:appendix-global}. Combining these identities proves the stated
necessary lower bound for every universal coefficient.

There is also a direct facet certificate for
\eqref{eq:appendix-invariant}, independent of the pyramid recursion.
Put $c=1/(d-1)!$. The facet area-normal/support pairs are
\[
 (-ce_1,0),\quad(-ce_2,0),\quad(-2ce_i,0)\quad(3\le i\le d),
\]
\[
 \left(c\left(e_1+\sum_{i=3}^d e_i\right),c\right),\qquad
 \left(c\left(e_2+\sum_{i=3}^d e_i\right),c\right).
\]
The corresponding halfspaces describe the square cross-sections
\[
 0\le x_1,x_2\le1-\sum_{i=3}^d x_i.
\]
Its volume is $2/d!$. Direct maximal-minor sums give
\[
 H=c^d d2^{d-1},\qquad L=c^{d+1}2^d,
\]
so $a=L/(d|C_d|H)=1/d$. To count $L$, only the four subsets obtained by
omitting either top facet or either of the first two coordinate facets
contribute, each by $c^{d+1}2^{d-2}$. For $H$, the contributions with zero,
one, or two top facets are respectively $1,d,d-1$, all multiplied by
$c^d2^{d-2}$.

The lower bound from the supplied small-truncation family is
\[
 L_d^{\mathrm{small}}=(d+1)
 \left[\frac{(d+1)^2}{d(d-1)}\right]^{1/(d-1)}.
\]
The square-pyramid bound is strictly larger since
\[
 \left(\frac{L_d^{\mathrm{square}}}{L_d^{\mathrm{small}}}\right)^{d-1}
 =\frac{d^2(d-1)}{d+1}>1\qquad(d\ge3).
\]
For example, its values for $d=3,4,5$ are approximately $13.8564$,
$13.5721$, and $14.0421$, respectively. Numerical values are illustrative
only; the formula is exact.

\begin{thebibliography}{9}
\bibitem{sharp}
\emph{Sharp lower-end simplex stability from integrated witnesses and
projection caps}. AI-assisted research supplement to entry005,
7 October 2026. Original sharp manuscript, pinned source:
\url{https://github.com/mxym/math/blob/e4352dd1a3415ee0dbe3213e13932173d879fe10/notes/sharp-simplex-stability/proof.tex}.

\bibitem{entry}
mxym, \emph{Random-determinant rigidity, sharp symmetric cone bounds, and
spectral nonattainment}. AI-assisted research manuscript, entry005,
version~3, 7 October 2026. Sections~1 and~4 give the invariant definition,
facet formula, and cone-law representation. Pinned source:
\url{https://github.com/mxym/math/blob/31e3d8a37e4a3051a7f5a2535be1c75642e15bcc/preprints/005-simplex-product-optimum/v3/paper.md}.

\bibitem{released}
mxym, \emph{Effective simplex rigidity for the projection cone invariant}.
Supplement to entry005, 7 October 2026. Pinned source:
\url{https://github.com/mxym/math/blob/87868bcb65bf0460a90d6bc2efc3324e46481230/notes/quantitative-projection-simplex-stability/paper.md}.

\bibitem{truncation}
\emph{An exact exponent obstruction for lower-end projection-cone simplex
stability}. Research supplement to entry005; AI-assisted research note,
7 October 2026. Pinned source:
\url{https://github.com/mxym/math/blob/87868bcb65bf0460a90d6bc2efc3324e46481230/notes/simplex-truncation-stability/proof.tex}.

\bibitem{classical}
R. Schneider, \emph{Convex Bodies: The Brunn--Minkowski Theory}, second
expanded edition, Cambridge University Press, 2014. Classical Cauchy,
surface-area, first-variation, and Brunn--Minkowski inputs; reference
inherited from the pinned source manuscripts.

\bibitem{boroczky}
K. B\"or\"oczky, Jr., \emph{The stability of the Rogers--Shephard inequality
and of some related inequalities}. Author-hosted manuscript:
\url{https://www.renyi.hu/~carlos/rogerstab.pdf}.

\bibitem{bh}
K. J. B\"or\"oczky and M. Henk, \emph{Cone-volume measure and stability}.
Author-hosted manuscript:
\url{https://www.renyi.hu/~carlos/cone-volume-stability.pdf}.

\bibitem{bfh}
K. J. B\"or\"oczky, F. Fodor and D. Hug, \emph{Strengthened inequalities
for the mean width and the $\ell$-norm}.
\url{https://arxiv.org/abs/2001.10706}.
\end{thebibliography}
\end{document}
